% ============================================================================
\section{Conclusion}
% ============================================================================

Nous avons construit une chaîne complète et reproductible : des workloads Go
compilés une fois et déployés sur quatre VM Google Cloud, un orchestrateur qui
exécute des files de tâches en parallèle et mesure chaque tâche, et un scheduler
qui décide où exécuter chaque tâche à partir d'un modèle calibré sur nos mesures.

Les mesures ont guidé l'optimisation. Elles ont montré que les VM ont la même
vitesse par tâche, qu'une VM à 2~vCPU ne vaut qu'environ 1,35 cœur, et que le coût
dépend avant tout des machines allumées. Notre algorithme \sel{} en tire les
conséquences : il choisit exactement combien de VM de chaque type allumer, puis
répartit les tâches avec Max-Min.

En conditions réelles, il réduit le coût de 71\,\% et 29\,\% sur les petits lots et
le temps de 24\,\% sur le gros lot, par rapport à Round Robin. Le modèle prédit les
temps à environ $\pm 10\,\%$ quand les VM fonctionnent normalement. La principale
limite est la variabilité des VM à cœur partagé, qui peut ralentir une machine de
moitié : c'est elle qu'il faudrait traiter en priorité, par un rééquilibrage
dynamique ou une mesure des VM juste avant de planifier.
